\section{Local Search and Adversarial Search}
\fatsf{Das Problem:}\\
Bisher:
\begin{itemize}
    \item Methoden, die den Suchraum systematisch durchsuchen, möglicherweise unter Verwendung von gezieltem \fatsf{Pruning} (z.B. A*)
    \item \enquote{Systematische Suchverfahren gehen \fatsf{alle Möglichkeiten geordnet durch} - und \fatsf{Pruning hilft, die schlechten Äste frühzeitig abzusägen}}
\end{itemize}
Aber:
\begin{itemize}
    \item Was passiert, wenn der Suchraum \fatsf{viel größer} wird?
    \item Was, wenn \fatsf{der Pfad nicht das Ziel ist}?
\end{itemize}
\subsection{Optimierungsprobleme}
Ein \fatsf{Optimierungsproblem} ist ein Problem, bei dem \fatsf{alle Zustände/Knoten als Lösung} betrachtet werden können (in unterschiedlichem Maße), aber das Ziel darin besteht, den Zustand zu finden, der eine Lösung nach einer \fatsf{Zielfunktion optimiert} (minimiert, maximiert, etc.). Das bedeutet, dass wir nach dem \fatsf{besten Zustand} suchen; es gibt keinen expliziten \fatsf{Zielzustand} und auch keinen \fatsf{Pfad (Kosten)}
\subsubsection*{Optimierungsproblem im Detail}
\begin{itemize}
    \item Eine \fatsf{Zielfunktion} (auch \fatsf{Evaluierungsfunktion}) gibt an, wie \enquote{gut} ein Zustand ist, und vergleicht ihn mit anderen Zustände. Der Wert der Zielfunktion kann je nach Problem entweder \fatsf{minimiert} oder \fatsf{maximiert} werden
\end{itemize}
\fatsf{Beispiel:}\\
Die \fatsf{Darwinistische Evolution} und das Überleben des Stärkeren können als \fatsf{Optimierungsprozess} betrachtet werden. Hier geht es nicht um ein einzelnes Ziel (z.B. den besten Individuum in einer Generation zu finden), sondern um die Evolution und Selektion im Laufe der Zeit
\clearpage
\subsubsection*{Wichtige Begriffe in Optimierungsproblemen}
\begin{enumerate}
    \item \fatsf{Konvergenz:} Dieser Begriff beschreibt die Eigenschaft einer Sequenz von Funktionswerten, die sich immer mehr einem bestimmten \fatsf{Grenzwert} oder \fatsf{Wert} annähert. In vielen Optimierungsprozessen strebt die Lösung nach einem optimalen Zustand
    \item \fatsf{Globales Optimum:} Ein \fatsf{globales Optimum} ist das Extremum (Minimum oder Maximum) der Zielfunktion für den gesamten Suchraum. Das bedeutet, es ist der bestmögliche Zustand, der über alle möglichen Zustände hinweg den besten Wert liefert
    \item \fatsf{Lokales Optimum:} Ein \fatsf{lokales Optimum} ist das Extremum der Zielfunktion für einen bestimmten Bereich des Suchraums. Zum Beispiel könnte es in einem Minimierungsproblem eine \enquote{Senke} in einem \fatsf{Talsystem} darstellen. Es ist der beste Zustand in diesem spezifischen Bereich, aber nicht unbedingt der beste Zustand im gesamten Suchraum 
\end{enumerate}
\begin{figure}[h]
    \centering
    \includegraphics[width=0.8\linewidth]{Local Search_Images/Key Terminologies.jpeg}
\end{figure}
\subsection{Optimierungsprobleme und Lokale Suchalgorithmen}
\fatsf{Optimierungsprobleme ohne Pfadabhängigkeit}\\
In vielen \fatsf{Optimierungsproblemen} sind die Reihenfolge der Aktionen und die damit verbundenen Kosten irrelevant.\\
Das Ziel besteht darin, einen Zustand zu finden, der alle \fatsf{Zielbedingungen} erfüllt - nicht den Pfad, um dorthin zu gelangen\\
\fatsf{Vorteile von Optimierungsansätzen:}
\begin{itemize}
    \item Benötigt nur eine \fatsf{geringe und konstante Menge an Speicher}
    \item Kann eine \fatsf{gute Lösung in sehr großen Suchräumen} finden, in denen systematische Methoden (z.B. A*) scheitern
\end{itemize}
\fatsf{Nachteile:}
\begin{itemize}
    \item Keine Garantie für \fatsf{Vollständigkeit} oder \fatsf{Optimallösung}
\end{itemize}

\subsection{Lokale Suchalgorithmen (Iterative Vierbesserungsmethoden)}
Lokale Suchalgorithmen sind Verfahren zur \fatsf{Lösung von Optimierungsproblemen}, bei denen \fatsf{nicht der gesamt Suchraum durchsucht wird}, sondern:
\begin{itemize}
    \item Es wird \fatsf{ein einzelner aktueller Zustand} betrachtet
    \item Man prüft dessen \fatsf{Nachbarn} (also ähnliche Zustände)
    \item Man wählt den \fatsf{besseren} Nachbarn und geht \fatsf{nur dorthin}
    \item Das wiederholt man \fatsf{so lange}, bis \fatsf{kein besserer Zustand} mehr gefunden wird
\end{itemize}
$\rightarrow$ Es geht also \fatsf{iterativ} und \fatsf{lokal} durch den Suchraum - immer nur auf Basis von kleinen Verbesserungen
\begin{figure}[h]
    \centering
    \includegraphics[width=0.8\linewidth]{Local Search_Images/Lokale Suche.jpeg}
\end{figure}
\subsubsection*{Wichtige Designüberlegungen für Suchalgorithmen}
Bei der Entwicklung eines lokalen Suchalgorithmus sind folgende Fragen entscheidend:
\begin{itemize}
    \item Wie wird das Problem dargestellt?
    \item Welche Zielfunktion soll optimiert werden?
    \item Wie wird die Qualität eines Zustands gemessen?
    \item Was definiert einen \enquote{Nachbarn} eines Zustands?
    \item Welche Zustände sind vom aktuellen Zustand aus erreichbar?
    \item Wie wird ein Nachbarzustand berechnet?
    \item Gibt es spezielle Einschränkungen, die ausgenutzt werden können?
\end{itemize}
\subsubsection*{Beispiele für Lokale Suche}
\begin{enumerate}
    \item \fatsf{Das Problem des Handlungsreisenden (Travelling Salesman Problem, TSP)}\\
    \fatsf{Problemstellung:}
    \begin{itemize}
        \item Ein Handlungsreisender muss eine Reihe von Städten besuchen und dabei die kürzeste Rundreise finden
    \end{itemize}
    \item \fatsf{Ansatz mit lokaler Suche:}
    \begin{itemize}
        \item Starte mit einer \fatsf{vollständigen, aber vermutlich suboptimalen Route}
        \item \fatsf{Verbessere iterativ} die Route durch \fatsf{paarweise Vertauschungen von Städten}
        \item Varianten dieser Methode können in großen Probleminstanzen sehr schnell eine \fatsf{Lösung finden, die weniger als 1\% von der optimalen Lösung abweicht}
    \end{itemize}
    \item \fatsf{Das n-Damen-Problem (n-Queens Problem)}\\
    \fatsf{Problemstellung:}
    \begin{itemize}
        \item Platziere \fatsf{n Damen} auf einem Schachbrett der Größe $n x n$, sodass \fatsf{keine Dame eine andere angreifen kann} (weder horizontal, vertikal noch diagonal)
    \end{itemize}
    \fatsf{Lokale Suche für das n-Damen-Problem}
    \begin{itemize}
        \item \fatsf{Zustände:} Eine Platzierung von n Damen in n Spalten (z.B. 256 mögliche Zustände bei n = 4)
        \item \fatsf{Nachbarschaftsoperationen:} Eine Dame in ihrer Spalte verschieben
        \item \fatsf{Zielfunktion:} h(n) = Anzahl der Konflikte zwischen Damen
        \item \fatsf{Zieltest:} Eine Lösung ist gefunden, wenn h(G) = 0 (keine Konflikte mehr)
    \end{itemize}
    \fatsf{Warum lokale Suche?}
    \begin{itemize}
        \item In jedem Schritt werden \fatsf{nur lokale Änderungen vorgenommen}
        \item Es wird sichergestellt, dass eine \fatsf{Folge lokaler Änderungen alle möglichen Zustände erreichen kann}
    \end{itemize}
\end{enumerate}
\begin{figure}[h]
    \centering
    \includegraphics[width=0.8\linewidth]{Local Search_Images/n-Queens Problem.jpeg}
\end{figure}
\subsection{Hill Climbing (Gradientenaufstieg) / Greedy Local Search}
\fatsf{Algorithmus:}
\begin{enumerate}
    \item \fatsf{Generiere alle Nachbarzustände} des aktuellen Zustands
    \item \fatsf{Wähle den Nachbarn mit der höchsten Bewertungsfunktion} (beste Heuristikwert)
    \item Wiederhole den Vorgang, \fatsf{bis keine Verbesserung mehr möglich ist} (das heißt, bis ein Maximum erreicht wird)
\end{enumerate}
\begin{figure}[h]
    \centering
    \includegraphics[width=0.8\linewidth]{Local Search_Images/Hill-Climbing.jpeg}
\end{figure}
\subsubsection*{Problem: Lokale Optima}
\fatsf{Herausforderung:}
\begin{itemize}
    \item Der Algorithmus \fatsf{stoppt, sobald er ein lokales Maximum erreicht}, selbst wenn dieses \fatsf{nicht das globle Optimum} ist
    \item Hindernisse wie \fatsf{Plateaus, Grate oder Schultern} können verhindern, dass die beste Lösung gefunden wird
    \item Dies wird oft verglichen mit: \enquote{Den Mount Everest im dichten Nebel mit Amnesie besteigen}
\end{itemize}
\subsubsection*{Lösungen: Varianten des Randomisierten Hill Climbing}
\begin{enumerate}
    \item \fatsf{Randomized Restart Hill Climbing}\\
    \fatsf{Idee:}
    \begin{itemize}
        \item \fatsf{Starte den Algorithmus mehrmals mit unterschiedlichen Anfangszuständen}, um verschiedene lokale Optima zu erreichen
        \item Durch mehrere Durchläufe steigt die Wahrscheinlichkeit, das \fatsf{globale Optimum} zu finden
    \end{itemize}
    \item \fatsf{Stochastic Hill Climbing}\\
    \fatsf{Idee:}
    \begin{itemize}
        \item \fatsf{Wähle den nächsten Zustand zufällig}, anstatt immer den besten Nachbarn zu nehmen
        \item \fatsf{Bessere Zustände haben eine höhere Wahrscheinlichkeit}, ausgewählt zu werden
        \item Dies verhindert, dass der Algorithmus \fatsf{zu früh in einem lokalen Optimum stecken bleibt}
    \end{itemize}
\end{enumerate}
\subsubsection*{Weitere Herausforderungen in Optimierungsproblemen}
\fatsf{Das Ridge-Problem (Grat-Problem)}
\begin{itemize}
    \item Jeder Nachbarzustand scheint schlechter zu sein (bergab), obwohl ein besserer Zustand existiert
    \item Das Problem entsteht, wenn der \fatsf{Suchraum in eine ungünstige Richtung geneigt} ist
    \item \fatsf{Metapher:} Ein gefaltetes Blatt Papier, das in einem unpassenden Winkel gehalten wird - jeder einzelne Schritt scheint bergab zu führen, obwohl das optimale Ziel bergauf leigt
\end{itemize}
\begin{figure}[h]
    \centering
    \includegraphics[width=0.3\linewidth]{Local Search_Images/Optimization Problems.jpeg}
\end{figure}
\subsubsection*{Beispiel: 8-Damen-Problem}
\fatsf{Herausforderung:}
\begin{itemize}
    \item Es gibt lokale Optima, bei denen \fatsf{keine Dame mehr bewegt werden kann}, ohne neue Konflikte zu erzeugen
    \item Beispielzustand:
    \begin{itemize}
        \item h = 17 $\rightarrow$ Startzustand
        \begin{figure}[h]
            \centering
            \includegraphics[width=0.3\linewidth]{Local Search_Images/h17.jpeg}
        \end{figure}
        \item \fatsf{Bester Nachbar:} h = 12 (Verbesserung)
        \item \fatsf{Lokal optimales h = 1 gefunden} $\rightarrow$ Keine Bewegung verbessert den Zustand mehr
        \begin{figure}[h]
            \centering
            \includegraphics[width=0.3\linewidth]{Local Search_Images/h1.jpeg}
        \end{figure}
        \item Was können Sie tun, um aus diesem lokalen Minimum herauszukommen?
    \end{itemize}
\end{itemize}
\subsubsection*{Mehrdimensionale Zustandslandschaften}
\fatsf{Suchlandschaft (State Landscape):}
\begin{itemize}
    \item Zustände werden als \fatsf{Positionen} dargestellt
    \item Der \fatsf{Heuristikwert} entspricht der \fatsf{Höhe} eines Geländes (hoch = gutes Ergebnis, tief = schlechtes Ergebnis)
\end{itemize}
\fatsf{Herausforderung:}
\begin{itemize}
    \item Je \fatsf{mehr Dimensionen}, desto \fatsf{schwieriger wird die Optimierung}
    \item \fatsf{Viele lokale Optima} erschweren Das Auffinden des globalen Optimums
\end{itemize}
\begin{figure}[h]
    \centering
    \includegraphics[width=1\linewidth]{Local Search_Images/State-Landscape.jpeg}
\end{figure}
\subsection{Gradientenabstieg (Gradient Descent)}
\begin{itemize}
    \item \fatsf{Gradientenabstieg ist eine zentrale Methode der Optimierung}, insbesondere für maschinelles Lernen und neuronale Netze
    \item Die Wahl der \fatsf{Lernrate} beeinflusst stark die Konvergenz
    \item \fatsf{SAT-Probleme} sind grundlegend für die theoretische Informatik und gehören zu den \fatsf{NP-vollständigen Problemen}
    \item \fatsf{Cook's Theorem} zeigt, dass SAT eine zentrale Rolle in der Komplexitätstheorie spielt
\end{itemize}
\subsection{Beam Search}
\fatsf{Idee:}
\begin{itemize}
    \item Statt nur \fatsf{einen Zustand} zu verfolgen, werden \fatsf{k Zustände} gleichzeitig gespeichert (\fatsf{k} = Strahlbreite)
\end{itemize}
\fatsf{Algorithmus:}
\begin{enumerate}
    \item Starte mit k zufällig generierten Zuständen
    \item Generiere alle Nachfolger aller k Zustände
    \item Wähle die k besten Nachfolger aus der gesamten Liste aus
    \item Wiederhole diesen Prozess, bis eine Lösung gefunden wird
\end{enumerate}
\fatsf{Vergleich:}
\begin{itemize}
    \item\fatsf{Hill-Climbing:} Speichert nur einen Zustand
    \item \fatsf{Beam Search:} Speichert mehrere Zustände (z.B. k = 2)
\end{itemize}
\begin{figure}[h]
    \centering
    \includegraphics[width=1\linewidth]{Local Search_Images/Hill-Climbing-Beam Search.jpeg}
\end{figure}
\subsection{Simulated Annealing (Simulierte Abkühlung)}
\fatsf{Idee:}
\begin{itemize}
    \item \fatsf{Kombiniert Hill-Climbing mit zufälligen Abwärtsbewegungen}, um lokale Optima zu vermeiden
    \item \fatsf{Zulassen von \enquote{schlechten} Schritten} (Bewegung in eine schlechtere Lösung), aber mit abnehmender Wahrscheinlichkeit über die Zeit
    \item Ähnelt einer \fatsf{Kombination aus Hill-Climbing und Random Walk}
\end{itemize}
\fatsf{Schlüsselelement: Temperatur (Hyperparameter)}
\begin{itemize}
    \item Die Temperatur T kontrolliert, wie oft \enquote{schlechte Züge} akzeptiert werden
    \item Die Temperatur \fatsf{nimmt exponentiell ab}, sodass mit der Zeit immer seltener schlechte Züge gemacht werden
\end{itemize}
\fatsf{Effektivität:}\\
Falls die Temperatur langsam genug gesenkt wird, \fatsf{konvergiert das Verfahren zum globalen Optimum}\\
\fatsf{Problem:} Kann extrem lange dauern
\subsection{Beispiel für Simulated Annealing}
\begin{figure}[h]
    \centering
    \includegraphics[width=1\linewidth]{Local Search_Images/Simulated Annealing.jpeg}
\end{figure}
\subsubsection*{Fazit: Simulated Annealing}
\fatsf{Grundprinzip:}
\begin{itemize}
    \item Simulated Annealing ist \fatsf{eine lokale Suche mit bewusst eingebautem Rauschen} das über die Zeit abnimmt
    \item Modelliert das Verhalten von \fatsf{physikalischen Systemen}, z.B. das Abkühlen von Metallen
\end{itemize}
\fatsf{Wissenschaftlicher Hintergrund:}
\begin{itemize}
    \item Verwendet Prinzipien aus der \fatsf{statistischen Physik} (Simulation von Gassystemen mit $10^{23}$ Teilchen)
    \item Entwickelt von \fatsf{Kirkpatrick et al. (1984)}
\end{itemize}
\fatsf{Allgemeine Anwendbarkeit:}
\begin{itemize}
    \item Kann genutzt werden, um komplexe Wahrscheinlichkeitsverteilungen zu simulieren
    \item Für \fatsf{Optimierung} wählt man eine geschickte Wahrscheinlichkeitsverteilung, die bei niedriger Temperatur auf das Optimum fokussiert ist
\end{itemize}
\subsection{Adversarial Search \& Spiele als KI-Herausforderung}
\subsubsection*{Warum sind Spiele für die KI-Forschung interessant?}
Spiele dienen als \fatsf{Herausforderung für KI-Algorithmen} und helfen Forschern, robuste maschinelle Lernmodelle zu entwickeln\\
\fatsf{Zitat:}\\
\enquote{Obwohl es interessant ist, dass DQN Spiele auf menschlichem Niveau spielt, könnte es diese Spiele auf eine ganz andere Weisen lernen als Menschen}\\
\fatsf{Warum Spiele?}
\begin{itemize}
    \item \fatsf{Messbarkeit:} Spiele ermöglichen es, den Fortschritt von KI-Algorithmen quantifizierbar zu verfolgen
    \item \fatsf{Sichere Trainingsumgebung:} Spiele sind simulationsfreundlich und haben geringes Risiko
    \item \fatsf{Öffentliches Interesse:} Größe KI-Meilensteine (z.B. AlphaGo gegen Lee Sedol) wecken Aufmerksamkeit und treiben Fotschritt voran
\end{itemize}
\subsubsection*{Adverarial Search - Was ist das?}
\fatsf{Definition:}\\
Adverasarial Search beschreibt Suchprobleme, bei denen \fatsf{mehrere Spieler mit entgegengesetzten Zielen} agieren\\
\fatsf{Eigenschaften:}
\begin{itemize}
    \item Spiele werden als \fatsf{Spielbäume} dargestellt
    \item Jeder Spieler muss die \fatsf{Aktion des Gegners antizipieren}
    \item Oft gibt es eine \fatsf{Zeitbegrenzung}, sodass \fatsf{keine vollständige Suche} möglich ist
\end{itemize}
\fatsf{Vergleich: Suchprobleme vs. Spiele}
\begin{table}[H]
    \centering
    \rowcolors{2}{\thepagecolor}{fgcolor!10!\thepagecolor}
    \renewcommand{\arraystretch}{1.2}
    \begin{tabular}{lll}
        \toprule
        Kriterium & Klassische Suche & Adversarial Search\\
        \midrule
        Gegner & Kein Gegner & Gegner mit eigenem Ziel\\
        Suchbaum & Agent sucht besten Weg zum Ziel & Spieler wählt besten Zug unter Berücksichtigung des Gegners\\
        Ziel & Fester Endzustand & Maximierung der eigenen Gewinnchance\\
        \bottomrule
    \end{tabular}
\end{table}
\subsubsection*{Adversarial Search: Wichtige Begriffe}
\begin{itemize}
    \item \fatsf{Perfekte vs. Imperfekte Information}
    \begin{itemize}
        \item Perfekte Information: Alle Spielzustände sind bekannt (z.B. Schach)
        \item Imperfekte Information: Teile des Spiels sind verborgen (z.B. Poker)
    \end{itemize}
    \item \fatsf{Deterministisch vs. Nicht-deterministisch}
    \begin{itemize}
        \item \fatsf{Deterministisch:} Die nächste Spielposition ist vollständig vorhersehbar (z.B. Schach)
        \item \fatsf{Nicht-deterministisch:} Ergebnisse haben Wahrscheinlichkeiten (z.B. Würfelspiele)
    \end{itemize}
    \item \fatsf{Zero-Sum-Spiele}
    \begin{itemize}
        \item Wenn ein Spieler gewinnt, verliert der andere exakt denselben Betrag
    \end{itemize}
\end{itemize}
\subsubsection*{Formalisierung eines Spiels als Suchproblem}
Ein Spiel kann als Suchproblem beschrieben werden mit:
\begin{enumerate}
    \item \fatsf{Initialzustand} $\rightarrow$ Anfangsposition des Spiels
    \item \fatsf{Spieler} $\rightarrow$ Gibt an, welcher Spieler am Zug ist
    \item \fatsf{Aktionen} $\rightarrow$ Alle möglichen legalen Züge eines Spielers
    \item \fatsf{Ergebnis} $\rightarrow$ Der resultierende Zustand nach einer Aktion
    \item \fatsf{Terminaltest} $\rightarrow$ Prüft, ob das Spiel endet
    \item \fatsf{Nutzenfunktion} $\rightarrow$ Gibt den Wert eines Endzustands für einen Spieler an
\end{enumerate}
\subsubsection*{Spielbäume \& ihre Struktur}
\begin{itemize}
    \item Ähnlich zu Suchbäumen, aber \fatsf{spielerspezifische Ebenen}
    \item Wurzelknoten = aktueller Spieler am Zug
    \item Blattknoten = Endzustände mit definiertem Nutzenwert
\end{itemize}
\fatsf{Beispiel: Tic-Tac-Toe}
\begin{itemize}
    \item Ein Tic-Tac-Toe Spielbaum hat \fatsf{weniger als 9! = 362880 Endzustände}
    \item Im Vergleich dazu: \fatsf{Schach hat $\approx$ $10^{40}$ mögliche Endzustände!}
\end{itemize}
\begin{figure}[h]
    \centering
    \includegraphics[width=0.8\linewidth]{Local Search_Images/Tic-Tac-Toe.jpeg}
\end{figure}
\subsubsection*{Herausforderungen bei der Adversial Search}
\begin{enumerate}
    \item \fatsf{Unvorhersehbarer Gegner}
    \begin{itemize}
        \item Der Algorithmus muss für \fatsf{jede mögliche Antwort des Gegners} eine Reaktion bereithalten
    \end{itemize}
    \item \fatsf{Zeitbegrenzung}
    \begin{itemize}
        \item In vielen Spielen kann \fatsf{nicht der gesamte Spielbaum durchsucht} werden $\rightarrow$ \fatsf{Heuristik \& Approximationen} nötig
    \end{itemize}
    \item \fatsf{Unterschied zwischen Spielen \& echten Problemen}
    \begin{table}[H]
        \centering
        \rowcolors{2}{\thepagecolor}{fgcolor!10!\thepagecolor}
        \renewcommand{\arraystretch}{1.2}
        \begin{tabular}{lll}
            \toprule
            Kategorie & Spiele & Reale Probleme\\
            \midrule
            Deterministisch? & Meistens ja & Oft nicht\\
            Gegenspieler? & Meistens nur einer & Oft mehrere Akteure\\
            Zugbasiert? & Ja & Meist nicht\\
            Ziel? & Gewinnen & Oft eine \fatsf{optimale Strategie}\\
            \bottomrule
        \end{tabular}
    \end{table}
\end{enumerate}
\subsection{Minimax \& Alpha-Beta Pruning}
\subsubsection*{Minimax-Algorithmus}
\fatsf{Grundidee}\\
Minimax nutzt \fatsf{Spielbäume}, um eine optimale Strategie für ein \fatsf{Zero-Sum-Spiel} zu finden. Dabei gibt es zwei Spielertypen:
\begin{itemize}
    \item \fatsf{MAX:} Versucht, den bestmöglichen Gewinn zu erzielen
    \item \fatsf{MIN:} Versucht, den Gewinn von MAX zu minimieren
\end{itemize}
\fatsf{Strategie:}
\begin{itemize}
    \item Spieler ziehen abwechseln
    \item MAX wählt die Aktion mit dem \fatsf{höchsten Minimax-Wert}
    \item MIN wählt die Aktion mit dem \fatsf{niedrigsten Minimax-Wert}
    \item \fatsf{Annahme:} Beide Spieler spielen \fatsf{optimal}
\end{itemize}
\fatsf{Eigenschaften des Minimax-Algorithmus}
\begin{itemize}
    \item \fatsf{Suchstrategie:} Tiefensuche (DFS)
    \item \fatsf{Komplettheit:} Ja, falls der Baum endlich ist
    \item \fatsf{Optimierung:} Optimal, wenn der Gegner optimal spielt
    \item \fatsf{Zeitkomplexität:} O($b^{m}$) (b = Verzweigungsfaktor, m = Tiefe)
    \item \fatsf{Raumkomplexität:} $O(bm)$ durch rekursive Berechnung
\end{itemize}
\fatsf{Problem:}
\begin{itemize}
    \item In vielen Spielen wächst der Spielbaum \fatsf{exponentiell}
    \item Beispiel: \fatsf{Schach} (b $\approx$ 35, m $\approx$ 100) $\rightarrow$ \fatsf{viel zu groß für vollständige Suche!}
\end{itemize}
\begin{figure}[h]
    \centering
    \includegraphics[width=0.6\linewidth]{Local Search_Images/Minimax Algorithm.jpeg}
\end{figure}
\subsubsection*{Heuristisch Bewertungsfunktion}
Das der vollständige Minimax-Baum oft zu groß ist, wird er \fatsf{abgeschnitten} und mit einer \fatsf{Heuristik} bewertet\\
\fatsf{Löusng:}
\begin{itemize}
    \item \fatsf{Begrenzte Suchtiefe} $\rightarrow$ Algorithmus sucht nur eine bestimmte Anzahl an Zügen voraus
    \item \fatsf{Bewertungsfunktion} $\rightarrow$ Schätzt, wie vorteilhaft eine Position ist
\end{itemize}
\fatsf{Beispiel für eine Heuristik in Schach:}
\begin{itemize}
    \item \fatsf{Materialvorteil} (z.B. Dame = 9 Punkte, Turn = 5 Punkte)
    \item \fatsf{Positionelle Faktoren} (z.B. Königsicherheit, Zentrumskontrolle)
\end{itemize}
\fatsf{Nachteil:}
\begin{itemize}
    \item Heuristiken sind nicht immer genau $\rightarrow$ \fatsf{Fehlereinschätzungen möglich!}
\end{itemize}
\subsubsection*{Alpha-Beta Pruning}
\fatsf{Grundidee}\\
Alpha-Beta Pruning \fatsf{optimiert Minimax}, indem es \fatsf{unnötige Berechnungen vermeidet} und trotzdem das gleiche optimale Ergebnis liefert\\
\fatsf{Prinzip:}
\begin{itemize}
    \item \fatsf{Alpha ($\alpha$):} Bestes bisher gefundenes Ergebnis für MAX (Start -$\infty$)
    \item \fatsf{Beta ($\beta$):} Bestes bisher gefundenes Ergebnis für MIN (Start +$\infty$)
\end{itemize}
\fatsf{Optimierung:}
\begin{itemize}
    \item Wenn \fatsf{$\alpha \geq \beta$}, dann \fatsf{wird der Teilbaum abgeschnitten (gepruned)}, da er nicht mehr relevant ist
    \item Dadurch werden \fatsf{keine unnötigen Knoten durchsucht}
\end{itemize}
\subsubsection*{Alpha-Beta-Prinning - Beispielablauf}
\begin{enumerate}
    \item \fatsf{Start bei MAX:}
    \begin{itemize}
        \item Setzt $\alpha = -\infty$, $\beta = +\infty$ und übergibt diese Werte an den ersten Kindknoten
    \end{itemize}
    \item \fatsf{MIN-Ebene:}
    \begin{itemize}
        \item Bestimmt das Minimum aus den Kindknoten und aktualisiert $\beta$
        \item Falls $\beta \leq \alpha$, wird der restliche Baum ignoriert
    \end{itemize}
    \item \fatsf{MAX-Ebene:}
    \begin{itemize}
        \item Bestimmt das Maximum aus den Kindknoten und aktualisiert $\alpha$
        \item Falls $\alpha \geq \beta$, wird der restliche Baum ignoriert
    \end{itemize}
\end{enumerate}
\begin{figure}[h]
    \centering
    \includegraphics[width=0.7\linewidth]{Local Search_Images/Alpha-Beta Pruning.jpeg}
\end{figure}
\fatsf{Ergebnis:}
\begin{itemize}
    \item \fatsf{Reduktion der Komplexität:}
    \begin{itemize}
        \item Minimax: O($b^{m}$)
        \item Alpha-Beta-Pruning: O($b^{m/2}$) $\rightarrow$ \fatsf{Halbierung der Suchzeit!}
    \end{itemize}
\end{itemize}
\fatsf{Optimierung:}
\begin{itemize}
    \item \fatsf{Best-First-Order:} Wenn die besten Züge zuerst geprüft werden maximiert sich der Effekt des Prunings
\end{itemize}
\subsubsection*{Grenzen von Alpha-Beta-Pruning}
\fatsf{Problem:}
\begin{itemize}
    \item Funktioniert schlecht in \fatsf{Spielen mit hohem Verzweigungsfaktor} (z.B. Go, b $\approx$ 250)
    \item \fatsf{Reine Suchalgorithmen reichen nicht aus}, um komplexe Spiele effizient zu meistern
\end{itemize}
\fatsf{Lösung: Monte-Carlo Tree Search (MCTS) \& KI-gestützte Evaluation}
\begin{itemize}
    \item \fatsf{AlphaGo (2016, DeepMind)} kombinierte \fatsf{Monte-Carlo-Tree-Search (MCTS) mit neuronalen Netzen}
    \item \fatsf{Lernfähige Bewertungsfunktionen} statt statischer Heuristiken
\end{itemize}
\subsubsection*{Warum ist Suche in der KI so zentral?}
\begin{itemize}
    \item \fatsf{Viele KI-Probleme sind schwer lösbar (interactable)}
    \item \fatsf{Suchalgorithmen sind oft die einzige Möglichkeit, komplexe Entscheidungen zu treffen}
    \item \fatsf{Beispiele für Such-basierte KI-Anwendungen:}
    \begin{itemize}
        \item \fatsf{Lernen:} Optimale Strategien durch Reinforcement Learning
        \item \fatsf{Planung:} Routenfindung für autonome Fahrzeuge
        \item \fatsf{Sprachverarbeitung:} Parsing natürlicher Sprache
        \item \fatsf{Bildverarbeitung:} Objekterkennung durch strukturierte Suche
    \end{itemize}
\end{itemize}
\fatsf{Fazit:} Obwohl Alpha-Beta-Pruning den Minmax-Algorithmus effizienter macht, sind für \fatsf{hochkomplexe Probleme moderner KI-Techniken}´wie \fatsf{MCTS \& neuronale Netze} notwendig